\label{chap:conclusions}

GraphMate ha sviluppato una pipeline riproducibile per rappresentare posizioni di
matto come grafi scacchistici e predire la prima mossa risolutiva. La progressione
dei modelli mostra che l'integrazione dei vincoli del dominio produce i miglioramenti
più rilevanti: il filtro di legalità migliora la baseline a vocabolario globale, lo
scoring contestuale delle mosse legali porta A3 al 67,56\% di Top-1 su Lichess e il
reranking delle posizioni successive alla mossa porta A4 all'85,41\% nel medesimo
benchmark.

\textbf{Risposta a RQ1.} Sul benchmark Lichess, il Modello B con timing supera
entrambe le baseline Qwen sottoposte al protocollo rigoroso UCI in tutte le
profondità disponibili, da mate-in-1 a mate-in-5. La conclusione riguarda il
confronto end-to-end congelato; non estende automaticamente il risultato ad altri
modelli linguistici, prompt o protocolli assistiti da strumenti.

\textbf{Risposta a RQ2.} Il timing sintetico implementato non migliora la
prestazione rispetto alla baseline A3 senza timing. La Top-1 di B è pari al
65,98\%, contro il 67,56\% di A3, con una differenza di (-1{,}58) punti
percentuali confermata dal confronto appaiato. Il risultato non dimostra che tempi
reali e allineati ai puzzle siano inutili; delimita l'esito alla codifica sintetica
valutata.

La valutazione esterna completa il quadro: A4 conserva un vantaggio descrittivo su
A3 nelle composizioni YACPDB, ma le accuratezze scendono sensibilmente rispetto a
Lichess e la differenza appaiata non è statisticamente significativa al 5\%. Il
cambio di distribuzione, la dimensione ridotta del benchmark, il limite Top-5 del
reranker, il timing sintetico e la predizione di una sola mossa sono i principali
vincoli interpretativi.

Nel complesso, gli esperimenti sostengono una conclusione precisa: per il compito e
i dati considerati, esplicitare legalità, candidate e conseguenze immediate delle
mosse è più utile dell'aggiunta del segnale temporale sintetico adottato. Allo
stesso tempo, il divario sul benchmark esterno mostra che accuratezza interna e
generalizzazione non coincidono e devono essere valutate separatamente.
